\documentclass[11pt]{article}
\usepackage[a4paper,margin=2.5cm]{geometry}
\usepackage[T1]{fontenc}
\usepackage{lmodern,microtype,amsmath,amssymb,booktabs,graphicx,adjustbox,xcolor,hyperref,natbib,amsfonts}
\hypersetup{colorlinks=true,linkcolor=blue!50!black,citecolor=blue!50!black,urlcolor=blue!50!black}
\setlength{\emergencystretch}{3em}  % long code names
\graphicspath{{figures/}}
\newcommand{\Hz}{H_0}
\newcommand{\Hv}{H_{\mathrm{vol}}}
\newcommand{\Hs}{H_\sigma}
\newcommand{\E}{\mathbb{E}}
\newcommand{\figureif}[3]{\IfFileExists{figures/#1}{\begin{figure}[htbp]\centering\includegraphics[width=#2\linewidth]{#1}\caption{#3}\label{fig:#1}\end{figure}}{}}
\newcommand{\tableif}[1]{\IfFileExists{tables/#1}{\input{tables/#1}}{}}
% Every number below that depends on data comes from tables/macros.tex (make report).
\IfFileExists{tables/macros.tex}{\input{tables/macros.tex}}{}
\providecommand{\resultHOne}{pending}
\providecommand{\resultHTwo}{pending}
\providecommand{\resultHThree}{pending}
\providecommand{\resultHOneBoth}{pending}
\providecommand{\resultHTwoBoth}{pending}
\providecommand{\resultHThreeBoth}{pending}
\providecommand{\venueBybitPairsPerp}{pending}
\providecommand{\venueBybitDiffPerp}{pending}
\providecommand{\venueBybitMaxDiffPerp}{pending}
\providecommand{\venueBybitPairsSpot}{pending}
\providecommand{\venueBybitDiffSpot}{pending}
\providecommand{\venueBybitMaxDiffSpot}{pending}
\providecommand{\venueOkxPairsPerp}{pending}
\providecommand{\venueOkxDiffPerp}{pending}
\providecommand{\venueOkxMaxDiffPerp}{pending}
\providecommand{\venueOkxPairsSpot}{pending}
\providecommand{\venueOkxDiffSpot}{pending}
\providecommand{\venueOkxMaxDiffSpot}{pending}
\providecommand{\nPairs}{60}
\providecommand{\nPerps}{59}
\providecommand{\pairDays}{pending}
\providecommand{\sdInflation}{1}
\providecommand{\resultFOne}{pending}
\providecommand{\resultFTwo}{pending}

\title{One Exponent to Rule Them All? An Asset-by-Asset Test of the Scaling Laws of
Order Flow, Volume, Volatility and Market Impact on Crypto Markets}
\author{nidhal.V}
\date{Working paper --- \today}

\begin{document}
\maketitle

\begin{abstract}
\citet{mocrs2026} derive from a nearly unstable Hawkes model of order flow that a single
statistic, the persistence $\Hz$ of the signed flow of core orders, fixes the roughness of
the volume rate ($\Hz-\tfrac12$), the roughness of extraday volatility ($2\Hz-\tfrac32$) and
the time exponent of the market impact of constant-rate meta-orders ($2-2\Hz$). Their
empirical analysis estimates $\Hz$ and the volume exponent on average over 40 stocks; the
volatility and impact exponents are taken from the literature. We estimate the four
exponents on the same data for each of \nPairs{} cryptocurrency pairs traded on Binance,
spot and perpetual (\nPerps{} perpetuals), with estimators validated by simulation under
microstructure noise, and test the three relations jointly in the cross-section with
errors-in-variables regressions and a bootstrap synchronized across assets, under a
pre-registered protocol with a frozen confirmation sample. Volume relation: \resultHOneBoth;
volatility relation: \resultHTwoBoth; impact relation (exploratory, proxy measure):
\resultHThreeBoth. Out of sample, an exact rough-volatility predictor is compared with HAR
benchmarks for daily volatility and volume, and the roughness implied by order flow with
the roughness measured on volatility (QLIKE ratios, spot, 22 days: \resultFOne{} and
\resultFTwo).
\end{abstract}

\section{Introduction}
Three facts of market microstructure are robust across assets and periods: volatility is
rough, with a Hurst exponent of log-volatility near 0.1 \citep{gatheral2018}; the signs of
market orders are positively correlated over thousands of trades
\citep{lillo2004,bouchaud2004}; and the price impact of a large order grows roughly as the
square root of its size \citep{toth2011}. \citet{eleuch2018} and \citet{jusselin2020}
connected them through nearly unstable Hawkes processes, and \citet{mocrs2026} unify them:
in their model a single exponent, $\Hz$, sets all three.

Their empirical support rests on equities and on averages: $\Hz$ estimated on 40 stocks
lies between 0.75 and 0.80 under a mixed fractional specification, and the roughness of
unsigned volume between 0.15 and 0.35, both consistent with $\Hz\approx 3/4$; the
roughness of volatility and the impact exponent are quoted from earlier studies. Whether
the relations hold \emph{asset by asset}, when the four exponents are measured on the same
data, is open. This paper tests it.

We contribute (i) estimates of the four exponents on the same trades for each asset;
(ii) a joint test of the predicted intercepts and slopes in the cross-section of assets,
with regressions that account for the estimation error of $\hat\Hz$ and a block bootstrap
synchronized across assets; (iii) a different market, open around the clock, where the
exchange reports the aggressor side of every trade; (iv) estimators validated beforehand on
simulated data with microstructure noise, which answers the objection that roughness can
be an artefact of estimation \citep{contdas2024}; and (v) a pre-registered protocol: the
hypotheses, estimators, grids and decision rule were fixed and tagged in the code
repository before any estimator ran on real data, and the confirmation sample was only
used once, with a frozen configuration. Four deviations, two of them decided after the
development sample, are listed with their reasons in Section~\ref{sec:deviations}.

A search of arXiv in November 2026, repeated every month since, found no test of these
relations on crypto markets; studies of Bitcoin report rough or
multifractal volatility \citep{takaishi2020} but do not relate it to the order flow.

\section{Theory and predictions}
\label{sec:theory}
\citet{mocrs2026} model \emph{core} buy and sell orders as two independent Hawkes
processes with a common kernel $\varphi_0$, and \emph{reaction} orders as a bivariate
Hawkes process driven by the core flow. In a nearly unstable regime the kernel norm tends to
one and its tail decays as a power law, $\int_t^\infty \varphi_0 \sim t^{-\alpha_0}$, with
$1/4<\alpha_0<1/2$. Their main results, in the scaling limit, are the following.

\begin{itemize}
\item \textbf{Signed flow.} The signed order flow converges to the sum of a process with
the local regularity of a fractional Brownian motion of Hurst index $\Hz=2\alpha_0\in
(1/2,1)$, coming from core orders, and of a martingale, coming from reaction orders; it is
well approximated by a mixed fractional Brownian motion $W_t + B^{\Hz}_t$ (their
Theorem~3.6 and Section~3.3).
\item \textbf{Volume.} Cumulative unsigned volume converges to the integral of a rough
process whose regularity is $\Hz-\tfrac12$ (Theorems~3.4--3.5): the prediction concerns
the \emph{volume rate}, not cumulative volume.
\item \textbf{Volatility and impact.} No-arbitrage arguments then give a volatility with
Hurst parameter $2\Hz-\tfrac32$, to be read as an \emph{extraday} volatility (their
Remark~4.2), and an average impact $\mathrm{MI}(t)\sim t^{2-2\Hz}$ in renormalised time
for a meta-order executed at a constant rate (their Theorem~4.1).
\end{itemize}

With $\Hz=3/4$, the square-root law of impact, the volume rate has roughness 1/4 and the
extraday volatility a Hurst parameter near 0. We write the three relations as
$Y_i=a+b\,H_{0,i}$ across assets $i$:
\begin{center}
\begin{tabular}{llll}
\toprule
Hypothesis & Relation & Prediction $(a;b)$ & Status \\
\midrule
H1 --- volume & $\Hv=\Hz-\tfrac12$ & $(-0.5;1)$ & confirmatory \\
H2 --- volatility & $\Hs=2\Hz-\tfrac32$ & $(-1.5;2)$ & confirmatory \\
H3 --- impact & $\delta=2-2\Hz$ & $(2;-2)$ & exploratory \\
H4 --- level & $\Hz\in[0.70,0.85]$ for liquid pairs & --- & descriptive \\
H5 --- robustness & same exponents on spot and perpetual & --- & robustness \\
\bottomrule
\end{tabular}
\end{center}
H3 is exploratory because without trader identifiers meta-orders can only be
approximated (Section~\ref{sec:estimators}). Two features of the source deserve mention:
one passage of its Section~3.3 reports mixed-model estimates around 0.65 while the
introduction and conclusion report 0.75--0.80, and the impact result is stated under the
approximations of its Section~4 without a proof in the appendix.

\section{Data}
\label{sec:data}
\paragraph{Source.} All data are the public files of Binance Vision: daily
\texttt{aggTrades} for spot pairs and USD-M perpetuals, and 1-minute klines. An aggTrade
groups the fills of one taker order at one price; the field \texttt{is\_buyer\_maker} gives
the side of the aggressor, so the sign $\varepsilon_i=+1$ for a buyer-initiated order and
$-1$ otherwise needs no classification rule. Each file is verified against its published
SHA-256 checksum. The data are licensed by Binance under CC BY-NC-SA 4.0 (Binance Vision
Dataset Terms, version 1.0); the statistics, tables and figures derived from them are
distributed under the same licence. This non-commercial research is not affiliated with or
endorsed by Binance. Other sources were reviewed: the Bybit archive has signed trades but no
checksums, and the other archives found were paid or incomplete.

\paragraph{Universe and samples.} USDT spot pairs, excluding stablecoin pairs, fiat and
leveraged tokens, are ranked by median daily quote volume over 2024-11-01 to 2026-10-31;
the first \nPairs{} with a daily aggTrades file on at least 95\,\% of days are kept, with the
perpetual of the same crypto when its files cover 95\,\% of days as well
(Table~\ref{tab:universe}). Availability is counted on aggTrades rather than klines:
klines can outlive trades. Symbol renames (MATIC to POL, RNDR to RENDER, FTM to S) are
merged through a versioned table; the days of a token swap are explained gaps. The gold
token PAXG and wrapped bitcoin (WBTC) were left out before any estimation: the first
follows the price of gold, the second duplicates bitcoin. The first twelve months form the
development sample, used to choose windows and ranges; the last twelve, the confirmation
sample, are estimated once with the frozen configuration. For daily volatility only, the
history extends to five years (2021-11 to 2026-10) where it exists, rebuilt from monthly
1-minute klines.
\tableif{universe.tex}

\paragraph{Quality controls.} A day enters the analysis only if it passes blocking controls:
checksum; schema; strictly increasing aggTrade identifiers; no missing aggTrade row, no run
of more than 1\,000 missing trade identifiers and at most 2\,\% of them missing; positive
prices and quantities; timestamps sorted inside the UTC day; and an OHLC rebuilt from the
trades that matches the 1-minute klines (no kline minute without trades, none off by more
than 1\,\%, at most 0.5\,\% of minutes off by more than 0.01\,\% or two ticks). Price moves
above 20\,\% within a second are flagged for inspection. Files re-published by Binance are
detected by their checksums and processed again. Two conventions of the perpetual feed
surfaced on the real files and are documented in the code repository: from mid-2025 a
small share of trade identifiers (0.1--0.3\,\%) belongs to fills that are not in the trade
files, while the aggTrade identifiers remain consecutive; and a 1-minute kline may open at the
last price of the previous minute. Neither indicates missing data, and the controls test
what does: missing rows, long runs of missing identifiers, missing minutes.
Table~\ref{tab:quality} reports the state of every pair-day of the window: \pairDays{}
pair-days are usable.
\tableif{quality.tex}

\paragraph{Reduction.} Each day is reduced, then its raw file is deleted. Trades are
aggregated on windows of $\Delta=10$\,s (and 1\,s for the ten most liquid pairs): signed
notional $F_k=\sum_{t_i\in\text{window }k}\varepsilon_i p_i q_i$, notional $V_k$, number of
trades and last price. Order flow has a strong weekly pattern; every series is divided by
its profile over the 168 hours of the week, estimated on the trailing 84 days only (no
look-ahead). From the windows we keep additive daily \emph{sufficient statistics}: sums
and sums of squares of the increments of $X_n=\sum_{k\le n}F_k$ for 30 log-spaced lags from
1\,s to 4\,h; sign autocovariances in event time up to $10^4$ trades, by FFT; sums of
$|v_{k+\ell}-v_k|^q$ for the volume rate $v_k=V_k/\Delta$; the 5-minute realised variance,
the two-scale realised variance \citep{zhang2005} and 96 local variances of 15 minutes;
and runs of same-side trades binned by notional. Every estimator and every bootstrap
replicate is a weighted sum of these daily rows, so no trade is read twice. On three random
days, each statistic equals a naive trade-by-trade computation to $10^{-14}$ in relative
terms; one day of the BTCUSDT perpetual during the liquidation cascade of 2025-10-10
(4.3 million aggTrades) is reduced in 14 seconds with 1.4\,GB of memory.

\section{Estimators}
\label{sec:estimators}
Each exponent has a principal and a control estimator. For $\Hz$ and $\delta$ a result is
reported as such only when the two agree within 0.05; for $\Hv$ and $\Hs$ the control is
shown but does not decide (deviation D2). For these two exponents the simulation study
chooses the principal among two candidates (Section~\ref{sec:simulation}).

\paragraph{$\Hz$.} The flow is first centred day by day in proportion to volume: each
window loses its volume times the day's imbalance (net flow over volume), so that a
day-level imbalance riding on bursts of activity is not read as memory (deviation D4;
in simulation, memoryless signs with a 5\,\% daily imbalance read as $\Hz\approx0.98$
without it). The principal estimator then fits the mixed model
\[
\E\bigl[(X_{t+\tau}-X_t)^2\bigr]=\sigma_1^2\,\tau^{2\Hz}+\sigma_2^2\,\tau
\]
by non-linear least squares on the aggregated variances, for $\tau$ from $\Delta$ to four
hours, with weights proportional to $1/\tau$ (the inverse sampling variance under the
Brownian null). Both terms take their exact expected shape after the centring, which
pins each day's path like a bridge. The fractional part must explain, at four hours,
more than twice the sampling error of the variance there and at least 5\,\% of it;
otherwise the flow is declared Brownian. The control is the quadratic-variation ratio on $(\Delta,2\Delta,4\Delta)$
used by \citet{mocrs2026}, in the spirit of the mixed-semimartingale estimators of
\citet{chong2025}; the simulation study preferred it to the event-time autocorrelation
of signs, $C(\ell)\sim\ell^{-\gamma}$ with $\Hz=1-\gamma/2$, which we report as a
diagnostic with the slope of a simple fBm (which reproduces the bias, increasing with
$\Delta$, noted by \citet{mocrs2026}) and a detrended fluctuation analysis of order two
\citep{peng1994}. The pre-registered estimator, on the uncentred flow, is a variant of
the specification curve. As a negative control, the signs are permuted within each day
and the same centring and fit must give $\Hz=1/2$ (deviation D3).

\paragraph{$\Hv$.} Two estimators on the deseasonalised volume rate. The first is that of
\citet{mocrs2026}: a generalised method of moments \citep{lixiu2016} on the autocovariance
of increments truncated at three standard deviations, which models white observation
noise explicitly. The second uses the moments
$m(q,\ell)=\overline{|v_{k+\ell}-v_k|^q}\propto\ell^{\,q\Hv}$ for
$q\in\{0.5,1,1.5,2,3\}$, with a check that $\zeta(q)$ is linear in $q$, on $v$ and on
$\log v$. The specification named the moments; the simulation study made the first
principal and the moments the control.

\paragraph{$\Hs$.} The prediction concerns extraday volatility, so the daily horizon is
primary, with two estimators on the daily 5-minute realised variance. The Whittle
estimator of \citet{fukasawa2022} models the realised variance as a day-averaged
fractional process observed with noise. The moments of increments of daily
log-volatility \citep{gatheral2018}, for lags of 1 to 50 days, were the estimator named by
the specification; after the simulation study the Whittle estimator is principal and the
moments are the control (deviation D1). The $p$-variation index of \citet{contdas2024}
and the intraday horizon (15-minute local variances) are reported; the five-year daily
history gives a robustness variant.

\paragraph{$\delta$.} Without account identifiers, meta-orders are approximated by runs of
same-side trades separated by less than $g$ seconds ($g\in\{1,5,30\}$): with
$I=\varepsilon(\log p_{\text{end}}-\log p_{\text{start}})$ and $Q$ the notional of the run,
we fit $\E[I\mid Q]=A\,Q^{\delta}$ on log-spaced bins of $Q$, and in the control the same law
against the duration of the run, which matches the time exponent of the theory. A run mixes
several traders, hence the exploratory status of H3. The aggregate impact of
\citet{patzelt2018} is reported as a description.

\section{Simulation study}
\label{sec:simulation}
Every estimator is first run on simulated data that go through the same reduction code as
the real files: exact fractional Gaussian noise by circulant embedding
\citep{daviesharte1987} and its mixture with a Brownian component for $\Hz$;
long-memory signs from the order-splitting model of \citet{lmf2005}; a rough volume rate
with observation noise for $\Hv$; rough Bergomi prices by the hybrid scheme of
\citet{bennedsen2017}, with microstructure noise and rounding to the tick, for $\Hs$; and
synthetic runs with a known impact exponent for $\delta$. The grid crosses true values
($\Hz\in\{0.6,0.7,0.75,0.8,0.9\}$, $\Hv\in\{0.1,0.2,0.25,0.3,0.4\}$,
$\Hs\in\{0.05,0.1,0.15,0.25\}$, $\delta\in\{0.3,0.4,0.5,0.6\}$) with noise-to-signal ratios
$\{0,0.5,1,2\}$; each cell has 200 replications of 365 days and 199 bootstrap draws.
An estimator is accepted where its bias is below 0.02 and the 95\,\% interval covers the
true value between 93\,\% and 97\,\% of the time.

The pre-registration fixes how the study decides: for $\Hv$ and $\Hs$ the estimator named
by the specification is principal if its bias is below 0.02 on the cells of the plausible
region, otherwise the alternative is, and the bootstrap spread is calibrated by the median
ratio of the true to the bootstrap standard deviation over those cells,
$k=\sdInflation$; coverage is then assessed on the calibrated intervals. A first, smaller
study had shown why: under long memory the block bootstrap underestimates the sampling
spread on short samples. The full study settled the choices: the autocovariance GMM is
principal for $\Hv$ (worst bias 0.014 against 0.079 for the moments); the quadratic-variation
ratio replaces the sign autocorrelation as control of $\Hz$ (whose bias reaches 0.033); for
$\Hs$ neither GJR ($+0.121$) nor Whittle (0.030) stays within 0.02 when
$\Hs\in[0.05,0.15]$: the rule kept GJR, a documented deviation makes the less biased
Whittle estimator principal (Section~\ref{sec:deviations}), and the volatility relation is
reported as limited by estimator reliability; $k=\sdInflation$.
Table~\ref{tab:simulation} and Figure~\ref{fig:fig5_bias_maps.pdf}
report bias, RMSE and coverage for every estimator, including the cells where it fails.
\tableif{simulation.tex}
\figureif{fig5_bias_maps.pdf}{0.95}{Bias of the principal and control estimators in
simulation; cell text: bias and coverage of the 95\,\% bootstrap interval.}

\section{Cross-sectional test}
\label{sec:inference}
\paragraph{Uncertainty.} Days are resampled by a stationary block bootstrap
\citep{politis1994} with a mean block of seven days (three and fourteen as sensitivity),
$B=999$, and the \emph{same} days for every pair and both markets: the strong correlation
between cryptocurrencies is kept in every replicate, and each replicate re-estimates every
exponent from the daily statistics. Interval half-widths are multiplied by $k$ and the
covariance of the test by $k^2$; results with $k=1$ are reported as well.

\paragraph{Regression and test.} $\hat\Hz$ is itself estimated, so least squares of $Y_i$
on $\hat H_{0,i}$ would be biased towards zero \citep{fuller1987}. We fit a Deming regression
with the ratio of error variances taken from the bootstrap, and test $(a,b)=(a_0,b_0)$
with the Wald statistic $W=(\hat\theta-\theta_0)^\top\hat\Sigma^{-1}(\hat\theta-\theta_0)
\sim\chi^2_2$, $\hat\Sigma$ being the bootstrap covariance of $(\hat a,\hat b)$; Holm's
correction \citep{holm1979} applies to H1--H3.

\paragraph{Decision.} A relation is \emph{rejected} if the corrected test rejects at 1\,\%
and the mean gap $\bar g=\overline{Y_i-a_0-b_0\hat H_{0,i}}$ exceeds 0.05 in Hurst units;
\emph{compatible} if the test does not reject and the 95\,\% interval of $\bar g$ lies
within $[-0.05,0.05]$; otherwise \emph{not decided}, with the number of days that would be
needed, assuming an interval width proportional to $1/\sqrt{T}$. Size and power are
measured on synthetic universes of assets that satisfy the relations, or violate them by
0.05, 0.10 and 0.20, with estimation errors of the size found by the bootstrap
(table~\ref{tab:power}).

\paragraph{Reading the decisions.} The joint test compares a fitted line with the
predicted one, so it needs a cross-section that lines up. Three situations leave a test
uninformative without making it wrong, and table~\ref{tab:reading} counts them. First,
the prediction may not be a value the exponent can take: $2\Hz-\tfrac32$ is negative as
soon as $\Hz<\tfrac34$, and no estimator of a Hurst exponent returns a negative value, so
H2 can only hold for pairs above that level. Second, an estimate may sit on a bound of
its search range and outside its own bootstrap interval; the interval of the mean gap is
then not centred on the gap, and two such estimates agree by construction. Third, when
the estimates of the two exponents are not correlated across pairs, the
errors-in-variables slope is unstable, its bootstrap covariance is large and the joint
test cannot reject whatever the mean gap: an undecided relation whose gap interval lies
far from zero is of this kind, and the gap is then the informative number.

\section{Protocol}
\label{sec:protocol}
The hypotheses, estimators, grids, decision rule and universe were written in
\path{config/preregistration.md} and committed with the tag \texttt{prereg-v1} before
any estimator ran on real data; until then real data served only to test ingestion and
quality. The configuration was frozen a second time, by the tag \texttt{config-frozen},
before the first estimation on the confirmation sample. Both rules are enforced by the
code, which refuses to estimate otherwise. The rule for choosing estimators from the
simulation study was amended before the full study had produced any result, in the light of
a preliminary small study: requiring raw bootstrap coverage would have retained estimators
known to be biased only because no estimator passes that test under long memory. Robustness checks --- windows $\Delta$,
ranges of $\tau$, without deseasonalisation, spot against perpetual, high and low
volatility regimes, without the five most liquid pairs, without days with extreme jumps,
block lengths, control estimators, five-year volatility --- are descriptive and shown
together as a specification curve.

\section{Results}
\label{sec:results}
\tableif{estimates.tex}
\figureif{fig1_h0_by_asset.pdf}{0.8}{$\hat\Hz$ by asset with 95\,\% bootstrap intervals.}
\tableif{tests.tex}
\tableif{reading.tex}
\tableif{levels.tex}
\figureif{fig2_h1_hvol.pdf}{0.85}{H1: volume roughness against $\hat\Hz$, with the
predicted line.}
\figureif{fig3_h2_hsigma.pdf}{0.85}{H2: volatility roughness against $\hat\Hz$, with the
predicted line.}
\figureif{fig4_h3_delta.pdf}{0.85}{H3: impact exponent against $\hat\Hz$, with the
predicted line.}
\figureif{fig6_forest_residuals.pdf}{0.95}{Per-asset gaps to the predicted relations with
95\,\% intervals.}

\section{Controls and robustness}
\label{sec:robustness}
\tableif{controls.tex}
\tableif{power.tex}
\figureif{fig7_spot_vs_perp.pdf}{0.95}{H5: exponents on the spot pair against the
perpetual of the same crypto.}
\figureif{fig8_specification_curve.pdf}{0.95}{Specification curve: every variant, sorted,
with its interval; not corrected for multiplicity.}

\section{Forecasting volatility and volume}
\label{sec:forecast}
The relations above are statements about exponents. A practical reading is predictive: if
volatility is rough with $H_\sigma = 2\Hz - \tfrac32$, a predictor that uses this roughness
should forecast daily volatility at least as well as the standard benchmarks, and the
roughness implied by order flow should forecast as well as the roughness measured on
volatility itself. We test both out of sample (pre-registration \S10); the evaluation is
statistical only, with no trading rule.

\paragraph{Design.} For each pair and market, at the end of every day $t$ of the evaluated
period, each model forecasts the average of the daily 5-minute realised variance (or of
the traded quote volume) over the next $h \in \{1, 5, 22\}$ days \citep{corsi2009}, from a
500-day window of the five-year daily history ending at $t$. Benchmarks are the random
walk, a direct AR(5) on logs and the HAR model on levels and on logs
\citep{corsi2009,abd2007}; HAR on logs is the reference. The rough models forecast the log
series by the Gaussian conditional expectation under the model of the Whittle estimator of
\citet{fukasawa2022}: the day average of a fractional Brownian motion plus white
measurement noise. The increments of day-averaged fBm have autocovariance
$\tfrac12\delta^4\Phi(k)$ with $\Phi(x) = |x|^{2H+2}/((2H+1)(2H+2))$, so the predictor is one
Toeplitz system per horizon, and its error variance gives the lognormal correction to
levels. Its parameters are estimated on the window (\emph{rough, $H$ estimated}), or $H$ is
fixed at the development estimate of $H_\sigma$, or at the order-flow value $2\hat\Hz - \tfrac32$.
The continuous-time kernel of \citet[\S5]{gatheral2018}, derived from \citet{nuzman2000},
is reported as published: on daily-sampled fBm it lost 13\,\% of one-step mean squared
error to the exact predictor in our simulations, because daily realised variance is an
average carrying measurement noise.

\paragraph{Evaluation.} Losses are QLIKE and the squared error on the level of the target,
which rank forecasts consistently with a noisy but unbiased volatility proxy
\citep{patton2011}. Per pair we report Diebold--Mariano tests against HAR
\citep{diebold1995,hln1997} with Newey--West variances \citep{neweywest1987} and the model
confidence set \citep{hansen2011}. The hypotheses pool pairs: the loss differential of each
pair is scaled by the reference's mean loss, averaged across pairs each day, and tested
with the same Diebold--Mariano statistic; a stationary bootstrap of dates gives the
interval of the average loss ratio. F1 compares the rough predictor with HAR, F2 the
order-flow $H$ with the volatility $H$ (the scaling law passes if they are equivalent within
2\,\%), F3 repeats F1 for volume. On simulated rough volatility the rough predictor lowers
QLIKE against HAR by about 1\,\%, 2--4\,\% and 7--13\,\% at 1, 5 and 22 days across
$H\in\{0.05,0.1,0.2\}$ with and without measurement noise, and wins in 70--100\,\% of the
replications (table~\ref{tab:forecast-sim}). Results: F1 \resultFOne; F2 \resultFTwo{}
(spot, 22 days).
\tableif{forecast.tex}
\tableif{forecast_sim.tex}
\figureif{fig9_forecast_ratios.pdf}{0.95}{Forecasting: mean QLIKE of the model relative to
the reference, pooled across pairs, with 95\,\% bootstrap intervals (blue: spot, orange:
perpetual).}

\section{Limitations}
The impact exponent rests on a proxy: runs of same-side trades mix independent traders, and
their duration is not the execution horizon of a meta-order. All assets trade on one
exchange; crypto order flow fragments across venues, and the persistence measured on one
venue need not equal that of the aggregate flow. An exploratory comparison with the
public trades of Bybit and OKX over one month (documentation, \emph{Other venues}) gives,
on perpetuals, a difference of $\Hz$ with Binance of \venueBybitDiffPerp{} on average for
Bybit (\venueBybitPairsPerp{} pairs, at most \venueBybitMaxDiffPerp{} in absolute value) and
of \venueOkxDiffPerp{} for OKX (\venueOkxPairsPerp{} pairs, at most \venueOkxMaxDiffPerp{});
on spot, \venueBybitDiffSpot{} and \venueOkxDiffSpot{}. On perpetuals these differences
are small next to the gap between spot and perpetuals measured on Binance (H5).
Twelve months is short for daily-horizon
roughness: on one year of daily data the estimate of $\Hs$ often sits on the lower bound
of its range (table~\ref{tab:reading}), which the five-year variant only partly addresses
since the history lacks the trades needed for $\Hz$. The perpetual trade files miss a small share of fills from mid-2025;
these fills appear to enter the klines but not the trade-level statistics. Finally, the relations
are tested under the assumption that each exponent is stable within a sample; the
specification curve and the comparison of the two samples show how far this holds.

\section{Deviations from protocol}
\label{sec:deviations}
Every change made after the tag \texttt{prereg-v1} is listed here with its reason and its
effect (\texttt{config/deviations.md}). D1 and D2 were decided on 2026-11-10 from the
simulation study alone, before any estimate on real data. D3 and D4 were decided on
2026-11-12, after the first estimates on the development sample and before any estimate on
the confirmation sample; the development results were then recomputed.

\textbf{D1: Whittle principal for $\Hs$.} The pre-registered rule keeps the default
estimator when neither candidate stays within 0.02 of the truth. For $\Hs$ that kept GJR,
whose bias reaches 0.121, over Whittle, whose bias reaches 0.030. The purpose of the rule is
to use the estimator the simulation finds reliable, so Whittle is principal and GJR the
control; H2 remains reported as limited.

\textbf{D2: concordance only against unbiased controls.} A result was to count only if the
principal and control estimators agree within 0.05. The controls of $\Hv$ (moments,
bias up to 0.079) and of $\Hs$ (GJR, 0.121) are biased by more than that margin, so a correct
result could fail the check because of the control. For these two exponents the control is
reported but does not decide whether the result counts; for $\Hz$ and $\delta$ the rule is
unchanged.

\textbf{D3: a negative control that can return 1/2.} Signs permuted within each day were
to give $\Hz=1/2$. On the development sample the control returned 0.80 on the flow and
1.00 on the autocorrelation of signs, for every pair. Neither could pass whatever the
truth: a power law fitted to a flat autocorrelation returns 1 by construction, and the
permutation keeps each day's imbalance, which rides on clustered activity and survives a
centring that is uniform in time. The permuted flow is now centred on each day's imbalance
in proportion to volume, like the principal estimate (D4), and the autocorrelation of
permuted signs is a diagnostic.

\textbf{D4: $\Hz$ on the flow centred day by day in proportion to volume.} The
pre-registered principal fitted the mixed model to the uncentred flow, which counts a
day-level imbalance as long memory: in simulation, signs without memory and a 5\,\% daily
imbalance read as $\Hz\approx0.98$. The principal, its variants, the control and the
diagnostics are now computed on the centred flow, with the exact expected variance of a
path pinned at both ends of the day; the pre-registered estimator remains a variant of the
specification curve. The $\Hz$ cells of the simulation study were run again and the
pre-registered acceptance rule applied: the worst bias in the plausible region is 0.0005,
and the calibration factor given by the rule moved from 1.14 to $k=\sdInflation$. The
choice rests on the definition of $\Hz$, the persistence of the flow and not the
day-to-day variation of its imbalance; it was not made on the tests, which it could move
either way. On the development sample it did: the volume relation on perpetuals went
from undecided to compatible, and the impact relation from undecided to rejected on both
markets; the other decisions did not change. Both estimators appear in the specification
curve.

\section{Conclusion}
One assumption links three exponents: if the persistence of order flow, $\Hz$, drives both
trading and prices, it sets the roughness of volume, the roughness of volatility and the
exponent of impact. On the confirmation year, with every choice frozen before that year
was read, the volume relation is \resultHOneBoth; the volatility relation is
\resultHTwoBoth; the exploratory impact relation, measured through a proxy, is
\resultHThreeBoth.

Every primary decision of the development year was reproduced on the confirmation year,
except the exploratory impact relation on perpetuals, which moved from rejected to
undecided. Two facts account for the pattern. The roughness of volume is the same on spot
and perpetuals (H5) while the persistence of order flow is higher on spot, so the volume
relation can hold on at most one of the two markets; it holds on perpetuals. And $\Hz$
lies below $3/4$ on both markets, where the volatility relation predicts a negative
exponent that no Hurst exponent can take; the daily estimates sit on the lower bound of
their range.

Two outcomes of the protocol hold whatever these decisions. The negative control, signs
permuted within each day, showed on the development year that the pre-registered
estimator counted a day-level imbalance as long memory. The corrected estimator (D4) was
checked in simulation and frozen with the rest of the configuration before the
confirmation year was used, so the confirmatory tests did not shape it. And the decision
rule keeps apart a relation that fails and one that cannot be judged: an interval wider
than the equivalence margin gives \emph{undecided}, not support.

Out of sample, the exact rough predictor forecasts the average volatility of the next 22
days on spot with a mean QLIKE ratio to HAR of \resultFOne{} (F1, below one when the rough
predictor is better); predicting with the roughness implied by order flow instead of the
roughness measured on volatility gives \resultFTwo{} (F2).

Three extensions would sharpen the tests: order flow aggregated across venues, since the
persistence on one exchange need not be that of the market; a longer trade history, so
that the roughness of daily volatility no longer sits on the bound of its range; and a
measure of impact from the order book, which the L2 collector of the pipeline records for
three pairs, in place of the proxy built on runs of same-side trades.

\appendix
\section{Reproducibility}
Every table and figure of this paper is rebuilt from a clean clone of the repository by
\texttt{make all}: download and verification of the data, reduction, simulation study,
estimation, tests and report. The pipeline also runs every day in a container, updates the
rolling estimates and publishes a dashboard and a read-only API. Code is released under
the MIT licence; raw Binance files are never redistributed, only derived statistics under
CC BY-NC-SA 4.0. Each estimate carries the version of the code and a hash of the
configuration it was computed with.

\bibliographystyle{plainnat}
\bibliography{references}
\end{document}
